% from ega2-1-fr.tex
\oldpage[II]{193}
Par abus de langage, on dira encore, en raison de l'existence du premier
isomorphisme (\hyperref[II.8.12.2.1-fr]{8.12.2.1}), que
$\mathcal{M}^{\square}$ est la \emph{fermeture projective} du
$\mathcal{O}_C$-Module $\widetilde{\mathcal{M}}$ (étant sous-entendu que dans
la donnée du $\mathcal{O}_C$-Module $\widetilde{\mathcal{M}}$ figure la
graduation du $\mathcal{S}$-Module $\mathcal{M}$).

\begin{env}[8.12.3]
\label{II.8.12.3-fr}
Avec les notations de (\hyperref[II.8.3.5-fr]{8.3.5}), on a un
homomorphisme canonique fonctoriel
\begin{equation}
\label{II.8.12.3.1-fr}
  p^*(\operatorname{Proj}_0(\mathcal{M}))
  \to\mathcal{M}^{\square}|\widehat{E}.
\tag{8.12.3.1}
\end{equation}
En effet, c'est un cas particulier de l'homomorphisme $u^{\sharp}$ défini de
façon générale dans (\hyperref[II.3.5.6-fr]{3.5.6}). Si
$Y=\operatorname{Spec}(A)$ est affine, $\mathcal{S}=\widetilde{S}$,
$\mathcal{M}=\widetilde{M}$, on voit, en se reportant à
(\hyperref[II.2.8.8-fr]{2.8.8}), que la restriction de
(\hyperref[II.8.12.3.1-fr]{8.12.3.1}) à
$p^{-1}(X_f)=\widehat{C}_f$ (pour un $f$ homogène dans $S_+$) correspond à
l'homomorphisme canonique
\begin{equation}
\label{II.8.12.3.2-fr}
  M_{(f)}\otimes_{S_{(f)}}S_f^{\leq}\to M_f^{\leq}
\tag{8.12.3.2}
\end{equation}
compte tenu de (\hyperref[II.8.2.3.2-fr]{8.2.3.2}) et
(\hyperref[II.8.2.5.2-fr]{8.2.5.2}).
\end{env}

\begin{env}[8.12.4]
\label{II.8.12.4-fr}
Plaçons-nous dans les hypothèses de (\hyperref[II.8.5.1-fr]{8.5.1}), dont
nous conservons les notations. Il résulte de
(\hyperref[II.1.5.6-fr]{1.5.6}) que pour tout $\mathcal{S}$-Module gradué
quasi-cohérent $\mathcal{M}$, on a d'une part un isomorphisme canonique
\begin{equation}
\label{II.8.12.4.1-fr}
  \Phi^*(\widetilde{\mathcal{M}})
  \xrightarrow{\sim}
  \left(q^*(\mathcal{M})\otimes_{q^*(\mathcal{S})}\mathcal{S}'\right)^{\sim}
\tag{8.12.4.1}
\end{equation}
de $\mathcal{O}_{C'}$-Modules~; et d'autre part
(\hyperref[II.3.5.6-fr]{3.5.6}) entraîne l'existence d'un
$\operatorname{Proj}(\varphi)$-morphisme canonique
\begin{equation}
\label{II.8.12.4.2-fr}
  \operatorname{Proj}_0(\mathcal{M})\to
  \left(\operatorname{Proj}_0
  (q^*(\mathcal{M})\otimes_{q^*(\mathcal{S})}\mathcal{S}')\right)
  |G(\varphi)
\tag{8.12.4.2}
\end{equation}
et aussi d'un $\widehat{\Phi}$-morphisme canonique
\begin{equation}
\label{II.8.12.4.3-fr}
  \operatorname{Proj}_0(\widehat{\mathcal{M}})\to
  \left(\operatorname{Proj}_0
  (q^*(\widehat{\mathcal{M}})
  \otimes_{q^*(\widehat{\mathcal{S}})}\widehat{\mathcal{S}'})\right)
  |G(\widehat{\varphi}).
\tag{8.12.4.3}
\end{equation}
\end{env}

\begin{env}[8.12.5]
\label{II.8.12.5-fr}
Considérons maintenant la situation de (\hyperref[II.8.6.1-fr]{8.6.1}),
avec les mêmes notations~; on prend donc dans ce qui précède $Y'=X$, le
morphisme $q:X\to Y$ étant le morphisme structural, et $\varphi$ est le
$q$-morphisme canonique (\hyperref[II.8.6.1.2-fr]{8.6.1.2}). On a alors un
isomorphisme canonique
\begin{equation}
\label{II.8.12.5.1-fr}
  q^*(\mathcal{M})\otimes_{q^*(\mathcal{S})}\mathcal{S}_X^{\geq}
  \xrightarrow{\sim}\mathcal{M}_X^{\geq}
\tag{8.12.5.1}
\end{equation}
en posant
$\mathcal{M}_X^{\geq}=\bigoplus_{n\geq0}\operatorname{Proj}_0(\mathcal{M}(n))$.
On peut en effet se borner au cas où $Y=\operatorname{Spec}(A)$ est affine,
$\mathcal{S}=\widetilde{S}$ et $\mathcal{M}=\widetilde{M}$, et définir
l'isomorphisme (\hyperref[II.8.12.5.1-fr]{8.12.5.1}) dans chacun des ouverts
affines $X_f$ ($f$ homogène dans $S_+$), en vérifiant la compatibilité avec le
passage à un multiple homogène de $f$. Or, la restriction à $X_f$ du premier
membre de (\hyperref[II.8.12.5.1-fr]{8.12.5.1}) est
$\widetilde{M}'=
\left((M\otimes_A S_{(f)})
\otimes_{S\otimes_A S_{(f)}}S_f^{\geq}\right)^{\sim}$ par
(\hyperref[II.8.6.2.1-fr]{8.6.2.1})~; comme on a un isomorphisme canonique de
$M\otimes_A S_{(f)}$ sur $M\otimes_S(S\otimes_A S_{(f)})$, on en déduit un
isomorphisme de $\widetilde{M}'$ sur
$\left(M\otimes_S S_f^{\geq}\right)^{\sim}$, et ce dernier est
canoniquement isomorphe par (\hyperref[II.8.2.9.1-fr]{8.2.9.1}) à la
restriction à $X_f$ du second membre de
(\hyperref[II.8.12.5.1-fr]{8.12.5.1}), et satisfait aux conditions de
compatibilité requises.

